\begin{afterword}
\summaryhead

No one can see the identity between quarks and a hand, but we can at least imagine atoms as billiard balls that push on things, as hands do. Anger seems different in kind. To a hunter-gatherer with no idea of computation, the gap between billiard balls and angry behaviour, including long-term plots, would look uncrossable, and a parody objector concludes that reductionism is flawed.

The author argues that talk of ``neurons'' that ``process information'', or the bare statement that anger is hormones, does not cross the gap. It is crossed by specific ideas: search trees, minimax and utility functions show how a merely causal system, such as a chess program, can compute plans, and a planner made of matter can be pushed by alcohol toward angry behaviour. No one can visualize the crossing from neurons to anger, as one can a hand made of fingers. Repeating words without understanding leads to the misreadings that anger is unjustified or unreal. The example is easy because it concerns behaviour, not experience, but it shows that explanatory gaps can be crossed with help from science.

\responsehead

The post's central claim is correct, and it is careful about its own scope. A system of parts that only push on each other can plan, if some of its internal states mirror the world and it searches among imagined actions for the one it rates highest. The idea has a history the post does not give: Kenneth Craik described minds that try out alternatives on a ``small-scale model'' of reality in 1943, and Claude Shannon proposed searching a chess game tree with minimax in 1950. The post does not claim the idea as its own. It limits itself to behaviour from the start, says so again at the end, and calls the case a practice problem.

That description is accurate, and it bears on who the post is answering. The objector is a parody, and the claim that non-materialists criticize this kind of materialism comes with ``I think'' and no one quoted. David Chalmers, whose zombie argument the book takes up later, counts explaining how a system can control behaviour among the ``easy problems'' and sees ``no obvious obstacle'' to explaining it in computational terms. Chalmers's objection concerns experience, which the post sets aside. So the gap the post crosses is one its opponents already grant can be crossed. The closing line, that ``Explanatory gaps can be crossed'', holds for this one. It does not apply to the gap Joseph Levine named, which was about experience.

Within its scope, the step to anger is the weakest part. The sketch shows how matter can plan. What makes a plan angry is left to a utility function the post does not describe, and the whole account of anger is one sentence in which alcohol pushes on ``the billiard balls making up the utility function.'' The post calls this ``tremendously oversimplified''.

In short: a correct sketch of how matter can produce planned behaviour, limited, as the post says, to behaviour, which is the part of the problem that the post's unnamed opponents do not dispute.
\end{afterword}
